\documentclass[10pt]{article}

% Specify the margins and text size.
\setlength{\textwidth}{6.4in}
\setlength{\textheight}{9.5in}
\setlength{\oddsidemargin}{0pt}
\setlength{\evensidemargin}{0pt}
\setlength{\topmargin}{0pt}
\setlength{\hoffset}{.05in}
\setlength{\voffset}{-.8in}

\setlength{\parskip}{5pt}
\setlength{\parindent}{0pt}

% Load some fonts and symbol packages
\usepackage{latexsym}
\usepackage{pifont}       % contains 'star' symbol for counterinsurgency handbook title
\usepackage{yfonts} 
\usepackage{amsmath}
\usepackage{amsfonts}

\usepackage{graphicx}     % actually, this is loaded by pstricks
\usepackage[T1]{fontenc}
\usepackage{ifthen}
\usepackage{pstricks,pst-func,pst-text,pst-node,multido,pst-plot,calc,pst-3dplot}
%\usepackage[all]{xy}
%\usepackage{animate}

% The hyperref package inserts links into the pdf.
\definecolor{MyLinkColor}{rgb}{.1,.2,1}
\definecolor{MyCiteColor}{rgb}{.1,1,.2}
\definecolor{MyURLColor}{rgb}{.4,.4,.4}
\usepackage[backref=true,pagebackref=false,hyperindex,colorlinks=true,
  linkcolor=MyLinkColor,urlcolor=MyURLColor]{hyperref}


% The tweaklist package is something I found on the web.  It provides a simple interface
% for making changes to spacing used in the itemize and enumerate environments.  Comment
% this out if you don't care to use tweaklist.
\usepackage{tweaklist}
\renewcommand{\itemhook}{\setlength{\parskip}{2pt}\setlength{\parsep}%
{1pt}\setlength{\topsep}{0pt}\setlength{\itemsep}{0pt}}

\newcommand{\N}{\ensuremath\mathbb{N}}
\newcommand{\Q}{\ensuremath\mathbb{Q}}
\newcommand{\R}{\ensuremath\mathbb{R}}
\newcommand{\C}{\ensuremath\mathbb{C}}
\newcommand{\z}{\hphantom{0}}

\begin{document}
\pagestyle{empty}

\href{http://www.su.edu}{\includegraphics[height=1.75cm]{sulogo.eps}}
\vspace{-1.69cm}

{\small{\ }\hfill
\begin{tabular}{cl}
& Math 207\\
& Introduction to Statistics\\
\end{tabular}
}
\medskip

{\setlength{\baselineskip}{1.05\baselineskip}

\rput(12,-1.2){\begin{pspicture}(0,0.5)(4,3.1)
\psset{yunit=0.75,dotsize=4pt}
\psaxes(5,4)
\psdot(1,2)\rput[bl](1.1,2.1){A}
\psdot(4,4)\rput[bl](4.1,4.1){B}
\psdot(5,3)\rput[bl](5.1,3.1){C}
\psdot(5,1)\rput[bl](5.1,1.1){D}
\psdot(3,0)\rput[bl](3.1,0.1){E}
\end{pspicture}}

1. For the figure at right, write down the\\ $x$ and 
$y$-coordinates of each point.
\vspace{0.5in}

2. Write down the equation of the line through  A and C.
\vspace{1in}

3. Write down the equation of the line through  A and E.
\vspace{1in}

4. Sketch the line $4\,y=x + 7$.
\vfill

\href{http://www.su.edu}{\includegraphics[height=1.75cm]{sulogo.eps}}
\vspace{-1.69cm}

{\small{\ }\hfill
\begin{tabular}{cl}
& Math 207\\
& Introduction to Statistics\\
\end{tabular}
}
\medskip


\rput(12,-1.2){\begin{pspicture}(0,0.5)(4,3.1)
\psset{yunit=0.75,dotsize=4pt}
\psaxes(5,4)
\psdot(1,2)\rput[bl](1.1,2.1){A}
\psdot(4,4)\rput[bl](4.1,4.1){B}
\psdot(5,3)\rput[bl](5.1,3.1){C}
\psdot(5,1)\rput[bl](5.1,1.1){D}
\psdot(3,0)\rput[bl](3.1,0.1){E}
\end{pspicture}}

1. For the figure at right, write down the\\ $x$ and 
$y$-coordinates of each point.
\vspace{0.5in}

2. Write down the equation of the line through  A and C.
\vspace{1in}

3. Write down the equation of the line through  A and E.
\vspace{1in}

4. Sketch the line $4\,y=x + 7$.}

\eject

5.  A large, representative sample of Americans was studied by the Public Health 
Service in the Health and Nutrition Examinaiton Survey (HANES2).  The
percentage of respondents who were left-handed decreased
steadily with age,
from 10\% at 20 years to 4\% at 70.  

\hspace{10pt} a) Sketch a possible histogram of the data.
\vspace{1in}


\hspace{10pt} b) \lq\lq The data show that many people change from left-handed to 
right-handed as they get older.\rq\rq\ True or false?  Why?  If false, how do you 
explain the pattern in the data?
\vspace{.5in}

6.  A personality test is administered to a large group of subjects.  
Five scores are shown below, in original units and standard units.  Fill in the blanks.\vspace{-5pt}
\newcommand{\bl}{\underbar{\hspace{20pt}}}
\begin{center}
{\setlength{\tabcolsep}{20pt}\begin{tabular}{rrrrr}
79 & 64 & 52 & 72 & \bl\\[15pt]
1.8 & 0.8 & \bl & \bl & $-1.4$
\end{tabular}}
\end{center}
\vfill

5.  A large, representative sample of Americans was studied by the Public Health 
Service in the Health and Nutrition Examinaiton Survey (HANES2).  The
percentage of respondents who were left-handed decreased
steadily with age,
from 10\% at 20 years to 4\% at 70.  

\hspace{10pt} a) Sketch a possible histogram of the data.
\vspace{1in}


\hspace{10pt} b) \lq\lq The data show that many people change from left-handed to 
right-handed as they get older.\rq\rq\ True or false?  Why?  If false, how do you 
explain the pattern in the data?
\vspace{.5in}

6.  A personality test is administered to a large group of subjects.  
Five scores are shown below, in original units and standard units.  Fill in the blanks.\vspace{-5pt}
\begin{center}
{\setlength{\tabcolsep}{20pt}\begin{tabular}{rrrrr}
79 & 64 & 52 & 72 & \bl\\[15pt]
1.8 & 0.8 & \bl & \bl & $-1.4$
\end{tabular}}
\end{center}
\vspace{.5in}

\end{document}
